\documentclass[11pt,a4paper]{article}

\usepackage[utf8]{inputenc}
\usepackage[T1]{fontenc}
\usepackage{lmodern}
\usepackage[margin=1in]{geometry}
\usepackage{hyperref}
\usepackage{xcolor}
\setlength{\parskip}{0.5em}
\setlength{\parindent}{0pt}

\hypersetup{
    colorlinks=true,
    linkcolor=blue!70!black,
    urlcolor=blue!70!black
}

\begin{document}

\begin{flushleft}
    \textbf{Quan Van} \\
    Independent Researcher, Ho Chi Minh City, Vietnam \\
    Email: \href{mailto:vanhaminhquan2406@gmail.com}{vanhaminhquan2406@gmail.com} \\
    ORCID: \href{https://orcid.org/0009-0002-8222-0076}{0009-0002-8222-0076} \\
    Date: October 15, 2026
\end{flushleft}

\vspace{0.5em}
\hrule
\vspace{1em}

\begin{flushleft}
    \textbf{To:} Editor-in-Chief and Editorial Board, \emph{Knowledge-Based Systems} \\
    \textbf{Subject:} Submission of Original Research Manuscript
\end{flushleft}

\vspace{0.5em}

Dear Editor-in-Chief and Editorial Board Members,

I am pleased to submit my original research manuscript entitled:

\begin{center}
\textbf{``Mining maximal high occupancy itemsets with a size-coupled upper bound''}
\end{center}

for consideration for publication as a regular research paper in \emph{Knowledge-Based Systems}.

\paragraph{Motivation and Core Contributions}
High-Occupancy Itemset (HOI) mining evaluates patterns that proportionally dominate the transactions containing them. However, since occupancy is not anti-monotone, HOIs can be overwhelmingly numerous, and existing miners (such as HEP and DFHOI) encounter severe scalability bottlenecks at low thresholds.

While condensed representations are well-studied for frequent patterns, mining \emph{Maximal HOIs (MHOIs)} directly has remained an open problem. Naive post-filtering of all HOIs suffers from combinatorial explosion, while filtering maximal frequent itemsets (MFIs) is mathematically incorrect. In this manuscript, we present:
\begin{enumerate}
    \item \textbf{Theoretical Foundations:} We prove that every MHOI is closed, demonstrate that an MHOI need not be an MFI, and establish that existing upper bounds (UBO) reduce trivially to support on databases with equal transaction lengths.
    \item \textbf{A Novel Size-Coupled Upper Bound ($B$):} We introduce a tight upper bound that constrains candidate extension sizes across supporting transactions, proving theoretically that it never exceeds UBO nor per-transaction bounds.
    \item \textbf{The MHOI-Miner Algorithm:} A depth-first search incorporating early child bounding, head-union-tail look-ahead, and subsumption pruning.
    \item \textbf{Rigorous Empirical Validation:} MHOI-Miner is verified against an exact oracle across 8,100 randomized instances. On standard benchmarks, it evaluates a median of 60.6$\times$ and 125$\times$ fewer itemsets than HEP+FILTER and DFHOI+FILTER, runs substantially faster in most settings, and successfully completes all benchmark configurations where existing miners exceed time or memory limits.
\end{enumerate}

\paragraph{Relevance to Knowledge-Based Systems}
Given that the state-of-the-art HOI miner, \textbf{DFHOI} (Nguyen et al., 2023), was published in \emph{Knowledge-Based Systems}, our paper directly advances this line of inquiry by proposing the first dedicated algorithm for maximal HOI discovery. We believe the theoretical insights and substantial performance gains will be of significant interest to your readership in pattern mining and knowledge discovery.

\paragraph{Declarations}
The author confirms that:
\begin{itemize}
    \item This manuscript is entirely original and has not been published or submitted elsewhere.
    \item All benchmark datasets are publicly accessible (SPMF, FIMI, and UCI repositories).
    \item There are no competing interests to declare.
\end{itemize}

Thank you very much for your time and consideration of this submission.

\vspace{1.5em}

Sincerely,

\vspace{0.8em}
\textbf{Quan Van} \\
\href{mailto:vanhaminhquan2406@gmail.com}{vanhaminhquan2406@gmail.com}

\end{document}
